\documentclass[conference]{IEEEtran}
\IEEEoverridecommandlockouts
% The preceding line is only needed to identify funding in the first footnote. If that is unneeded, please comment it out.

\usepackage{cite}
\usepackage{amsmath,amssymb,amsfonts}
\usepackage{algorithmic}
\usepackage{graphicx}
\usepackage{textcomp}
\usepackage{xcolor}
\usepackage{booktabs} % For formal tables
\usepackage{hyperref} % For clickable links

\def\BibTeX{{\rm B\kern-.05em{\sc i\kern-.025em b}\kern-.08em
    T\kern-.1667em\lower.7ex\hbox{E}\kern-.125emX}}

\begin{document}

\title{QREM Core: A Highly-Efficient Hardware Architecture for the ML-KEM PQC Standard\\
% {\footnotesize \textsuperscript{*}Note: Sub-titles are not captured in Xplore and
% should not be used}
}

\author{\IEEEauthorblockN{1\textsuperscript{st} Kiet Le}
\IEEEauthorblockA{\textit{Lassonde School of Engineering} \\
\textit{York University}\\
Toronto, Ontario, Canada \\
kietle@my.yorku.ca}
\and
\IEEEauthorblockN{2\textsuperscript{nd} Team Member 2}
\IEEEauthorblockA{\textit{Lassonde School of Engineering} \\
\textit{York University}\\
Toronto, Ontario, Canada \\
email@my.yorku.ca}
\and
\IEEEauthorblockN{3\textsuperscript{rd} Team Member 3}
\IEEEauthorblockA{\textit{Lassonde School of Engineering} \\
\textit{York University}\\
Toronto, Ontario, Canada \\
email@my.yorku.ca}
% Add more authors as needed
}

\maketitle

\begin{abstract}
The rapid advancement of quantum computing threatens the security of classical public-key cryptography. In response, NIST standardized the Module-Lattice-Based Key-Encapsulation Mechanism (ML-KEM) under FIPS 203. This paper presents the QREM Core, a highly efficient, cycle-accurate hardware accelerator for ML-KEM. Our architecture features a dedicated Polynomial Arithmetic Unit (PAU) optimized for 12x12-bit modular multiplication, a high-throughput Hash Sampler Unit (HSU) for FIPS 202 Keccak functions, and a deterministic 4-bank polynomial memory subsystem. By strategically decoupling fixed-point compression into a dedicated Transcoder Unit, we eliminate combinational area bloat, achieving a highly optimized ASIC and FPGA footprint. The fully integrated system is verified against a custom Python Golden Model and standard NIST Known Answer Tests (KATs).
\end{abstract}

\begin{IEEEkeywords}
Post-Quantum Cryptography (PQC), ML-KEM, FIPS 203, Hardware Accelerator, ASIC, FPGA, Number Theoretic Transform (NTT)
\end{IEEEkeywords}

\section{Introduction}
% Introduce the problem: Quantum threat, Shor's algorithm, NIST standardization.
% Introduce the solution: Hardware acceleration of ML-KEM.
% Briefly outline the contributions of this paper.

\section{Background}
\subsection{Module-Lattice-Based Cryptography}
% Mathematical foundation of ML-KEM (R_q, module dimension k).

\subsection{The FIPS 203 Standard}
% KeyGen, Encapsulation, Decapsulation processes.

\section{System Architecture}
% High-level block diagram explanation.
The QREM Core architecture is partitioned into several independent datapaths communicating via AXI4-Stream interfaces...

\subsection{Polynomial Arithmetic Unit (PAU)}
% NTT/INTT butterfly operations and modular reduction.

\subsection{Hash Sampler Unit (HSU)}
% SHA3-256, SHA3-512, SHAKE128/256 for seed expansion and CBD sampling.

\subsection{Transcoder Unit}
% ByteEncode/Decode and 12x24-bit fixed-point Compression math.

\subsection{Poly-Mem Subsystem}
% 4-bank SRAM memory routing, arbitration (PAU > HSU > Transcoder).

\section{Verification Methodology}
% Describe the TAID methodology, Python Golden Model, Verilator regression, and Yosys Scatter-Gather pipeline.

\section{Implementation Results}
\subsection{Area and Resource Utilization}
% FPGA (LUTs, DSPs, BRAMs) and ASIC (Gate Equivalents) results.

\subsection{Performance and Throughput}
% Latency cycles and F_max.

\section{Conclusion}
% Summarize the achievements and future work.

\section*{Acknowledgment}
% Thank supervisors, Lassonde Hardware Design Club members, or organizations that supported the capstone.

\bibliographystyle{IEEEtran}
\bibliography{references} % Requires a references.bib file

\end{document}
